\section{STaR：用推理自举推理}

\subsection{少量 rationale 如何扩展成大数据}

STaR（Self-Taught Reasoner）从少量带 reasoning step 的示例开始，让模型为大量只有问题--答案的数据生成 rationale；只保留最终答案正确的轨迹，加入训练集微调，再重复这一过程。

\videofigure{figures/star-bootstrapping.jpg}{STaR 从少量 reasoning examples 迭代生成、过滤和训练。}{00:17:30--00:18:42}

若训练集为 $\{(x_i,a_i)\}$，第 $m$ 轮模型生成 rationale $r_i$：

$$
r_i\sim p_{\theta_m}(r\mid x_i),\qquad
\mathcal D_m=\{(x_i,r_i,a_i):\hat a(r_i)=a_i\}.
$$

然后在累积的正确 rationale 上训练：

$$
\theta_{m+1}=\arg\max_\theta
\sum_{(x,r,a)\in\mathcal D_m}\log p_\theta(r,a\mid x).
$$

\subsection{Rationalization：从已知答案反向解释}

困难问题上，模型可能始终无法自行生成正确答案，因而没有新正例。STaR 的关键创新是把正确答案提供给模型，要求它生成一个能解释该答案的 rationale。

\videofigure{figures/rationalization.jpg}{Rationalization 用正确答案帮助模型为原本失败的问题构造解释。}{00:18:42--00:20:00}

\begin{knowledgebox}{Rationalization 类似 answer-conditioned proposal}
它不是证明模型原先会解题，而是利用答案作为额外条件搜索一条解释。若解释质量可靠，这能把困难问题加入训练；若模型只是事后编造看似合理的故事，就会引入伪推理。
\end{knowledgebox}

\subsection{完整自举循环}

\videofigure{figures/star-contributions.jpg}{STaR 通过生成、正确性过滤、rationalize 与再训练迭代扩展 rationale 数据。}{00:20:00--00:22:42}

\videofigure{figures/star-method.jpg}{Vanilla STaR 与加入 rationalization 的数据流。}{00:24:50--00:27:28}

算法可概括为：

\begin{enumerate}
    \item 用当前模型为所有问题生成 rationale 与答案；
    \item 按最终答案过滤正确轨迹；
    \item 对失败问题，在给定正确答案条件下 rationalize；
    \item 用新 rationale 数据微调模型；
    \item 重复直到收益饱和或数据质量下降。
\end{enumerate}

\subsection{隐藏假设与错误模式}

\videofigure{figures/star-assumptions.jpg}{STaR 假设最终正确性可以近似 reasoning 质量。}{00:22:42--00:24:50}

这个假设会产生：

\begin{itemize}
    \item \textbf{False positive rationale}：过程错误但最终答案碰巧正确；
    \item \textbf{False negative rationale}：过程有价值但最终小错，被整条丢弃；
    \item \textbf{Selection bias}：训练数据越来越集中在模型已经能解的问题；
    \item \textbf{Confirmation loop}：模型偏好某类解释，下一轮继续复制同一模式。
\end{itemize}

\subsection{结果与 Rationalization Gap}

\videofigure{figures/star-results.jpg}{STaR 在 commonsense 与数学 reasoning 上优于无 scratchpad 或仅答案微调。}{00:27:28--00:31:22}

\videofigure{figures/rationalization-challenge.jpg}{给定答案生成的 rationale 与无答案推理存在分布差异。}{00:31:22--00:33:34}

Rationalization 训练数据条件中含答案，但部署时没有答案；模型可能学习只有在答案已知时才成立的解释模式。这一 distribution mismatch 是 STaR 的核心局限。

\begin{warningbox}{最终答案过滤不是 process verification}
STaR 的 reward 只检查终局，因此不能识别中间逻辑漏洞。更稳健的扩展可以引入 PRM、formal checker 或多个 verifier，对 rationale 本身做质量控制。
\end{warningbox}

\subsection{本章小结}

STaR 展示了低成本 reasoning data bootstrapping：用模型生成自己的训练数据，并通过可验证答案过滤。它样本效率高、实现简单，但容易受终局标签和 answer-conditioned rationalization 偏差限制。

